\documentclass[11pt]{article}
\usepackage[a4paper,margin=24mm]{geometry}
\usepackage{amsmath,amssymb,amsthm,booktabs,microtype}
\usepackage[hidelinks]{hyperref}
\newcommand{\E}{\mathbb E}\newcommand{\Pp}{\mathbb P}
\newtheorem{theorem}{Theorem}\newtheorem{lemma}[theorem]{Lemma}
\title{A geometric region closed by the Bernoulli reserve}
\author{}\date{28 October 2026}
\begin{document}\maketitle

The geometric budget and the Bernoulli reserve have the same quadratic
factor $Q(p)$. We use this to close an entire part of the remaining left
range. The proof keeps all phases and all admissible noise in that region.
It uses no count-tail certificate: fourth moments give a summable bound
for the signed curvature in the exact crossing identity.

Write $\phi_0=(2\pi)^{-1/2}$, $\sigma=\sqrt{10}$ and
\[
 C_E=(3+\sigma)\phi_0/6,\quad
 Q(p)=(6+2\sigma)(p-p_*)^2,\quad p_*=(4-\sigma)/2.
\]
For $Y=B+U$, put $p=\Pp(B=1)$, $q=1-p$, $v=pq$,
$s=\E U^2$, $\lambda=ns$, $\ell=(v+s)^{-1/2}$ and $X=\ell(Y-p)$.
The two conditional laws $U_0,U_1$ are centered. The notation $D$ denotes
the actual support-gap budget of the v13 geometric theorem:
\[
 D=(3+\sigma)\E|X|^3-\E X^3-3\sum_{(a,b)}\tau_{a,b}(b-a)^2.
\]

\begin{theorem}[An excluded probability--index region]\label{g:main}
Let $n\ge42$ be an integer and $1/3\le p\le2/5$. Suppose
\[
 -7/10\le U_0\le a_0,\qquad -a_1\le U_1\le7/10,
 \quad a_0,a_1\ge0,\quad a_0+a_1\le17/250,
\]
\[
 0\le\lambda\le8/25,\qquad nD\le129/4,\qquad |H|\le46/25.
\]
If $nQ(p)/\sqrt v\ge24$, then, for the sum $S_n$ of $n$ iid copies of $X$,
\[
 \sqrt n\{\Pp(S_n\le H)-\Phi(H/\sqrt n)\}
                  -C_E\E|X|^3<-\frac1{10n}.
\]
\end{theorem}

The acquired budget $nD\le32.20665378$ is smaller than $129/4$;
the cap $\lambda<8/25$ on this probability range is already proved in
v13. Thus an acquired witness with nonnegative signed excess must satisfy
\begin{equation}\label{g:newindex}
                         nQ(p)/\sqrt{pq}<24.
\end{equation}
This improves the previous necessary bound $nQ(p)/\sqrt{pq}\le32.20665378$
by using the actual CDF, rather than by slightly changing a noise ceiling.
For $1/3\le p\le b$, this gives the following integer-index bounds:
\[
\begin{array}{c|ccccc}
 b&1/3&.35&.375&.39&.40\\\hline
 n\text{ at most}&125&195&490&1140&2681.
\end{array}
\]
Indeed $Q(p)/\sqrt{pq}$ decreases on this interval, and the exact endpoint
comparisons use $n=126,196,491,1141,2682$, respectively. The previous v13
bound $n\le168$ for $p<1/3$ extends every column with $b\ge.35$ to the
whole acquired interval $[.25,b]$. It does not extend the first column.
The finite base and the remaining region below \eqref{g:newindex} are not
settled. No quantitative Bernoulli reserve is asserted below $p=1/3$.

\section{A summable bound for the signed curvature}
Let $W_j$ be a sum of $n-j$ independent copies of $U_0$ and $j$ of $U_1$.
Set $\pi_j=\Pp(\operatorname{Bin}(n,p)=j)$ and
\[
 B_j(t)=\E(t-j-W_j)_+-(t-j)_+,\qquad
 \mathcal C(t)=\sum_{j=1}^{n-1}(\pi_{j+1}-\pi_j)B_j(t).
\]
The following lemma is independent of the geometric budget and of the
Bernoulli theorem. It is useful on either probability range.

\begin{lemma}\label{g:curvature}
Let $n\ge2$, $0<p<1$, and put $c=(np-q-t)_+\le1$.
Suppose $T,M\ge0$ satisfy $\E W_j^2\le T$ and $\E W_j^4\le M$
for every $j$ with $\pi_{j+1}>\pi_j$. If
\[
 \max_r\Pp(\operatorname{Bin}(n+1,p)=r)
                  \le\frac{\beta}{\sqrt{(n+1)v}},
\]
then
\begin{equation}\label{g:curvbound}
 \sqrt n\,\mathcal C(t)\le\frac{\beta}{nv^{3/2}}
 \left\{\frac T4+c\sqrt T+
                 \frac{27M}{256}\left(\frac53+\frac{5c}4\right)\right\}.
\end{equation}
For bounded conditional laws $|U_j|\le A$, one may take
$T=\lambda/q$ and $M=A^2T+3T^2$.
\end{lemma}
\begin{proof}
For $a=t-j$, centering gives
\[
 B_j=\tfrac12\{\E|a-W_j|-|a|\}\ge0,
 \qquad B_j\le\sqrt T/2,\qquad a_+B_j\le T/4.
\]
The last bound follows from
$\sqrt{a^2+T}-a\le T/(2a)$ for $a>0$. Also
$(x-a)_+\le27x^4/(256a^3)$ for $a>0$: the maximum of
$(x-a)/x^4$ is attained at $x=4a/3$. Consequently
\[
                       B_j\le27M/(256|a|^3)\quad(a\ne0).
\]
The exact adjacent-mass identity, with
$\widetilde\pi_r=\Pp(\operatorname{Bin}(n+1,p)=r)$, is
\[
 \pi_{j+1}-\pi_j
    =\frac{np-q-j}{(n+1)v}\widetilde\pi_{j+1}.
\]
Discard only its nonpositive terms. Its positive numerator is at most
$(a+c)_+$. Among the remaining integer mesh points there is at most one
with $-c<a<0$ and at most one with $0\le a<1$. Their total is bounded
by $c\sqrt T+T/4$. For $a\ge1$, list the points as $a=b+i$, $i\ge1$,
with $0\le b<1$. The fourth-moment bound and
\[
 \sum_{i\ge1}i^{-2}<5/3,\qquad \sum_{i\ge1}i^{-3}<5/4
\]
give the last term in braces. Both series bounds follow from their first
few terms and an integral remainder. Multiply by the mass bound and use
$\sqrt n/(n+1)^{3/2}\le1/n$.

Finally, a positive mass difference requires $j<np-q<np$. Writing
$s_i=\E U_i^2$, the variance
$(n-j)s_0+js_1$ is at most $\lambda/q$, because its coefficients relative
to $nqs_0$ and $nps_1$ are at most $1/q$ and $1$, respectively.
Independence gives
$\E W_j^4=\sum\E U_i^4+3(\E W_j^2)^2-3\sum(\E U_i^2)^2
\le A^2T+3T^2$. The moment and count allocations remain those of the
same two conditional laws.
\end{proof}

For $p\in[1/3,2/5]$ and $n\ge42$ one may take $\beta=41/100$.
Here is a short finite justification. Put $N=n+1\ge43$, and let $r=m/N$
for a binomial mode $m=\lfloor(N+1)p\rfloor$. Then
$-q\le m-Np\le p$ and
\[
 r(1-r)\ge v\{1-1/N-2/N^2\}>39v/40.
\]
Finite Stirling, after discarding the nonpositive entropy exponent and
negative denominator remainders, bounds the maximum mass by
\[
 \frac{\phi_0}{\sqrt{Nr(1-r)}}e^{1/(12N)}
 <\frac{(2/5)\sqrt{40/39}(516/515)}{\sqrt{Nv}}
 <\frac{41/100}{\sqrt{Nv}}.
\]
The final comparison is checked by squaring positive rational quantities.

\section{The geometric cost leaves less than 0.09 noise}
The v13 inward-cubic inequality, with $\bar d=17/250$, gives
\[
 nD\ge\frac{nQ/\sqrt v+
           \lambda\min(Y_0,Y_1)/v^{3/2}}
                {(1+\lambda/(nv))^{3/2}},
\]
\[
 Y_0=3(\sigma+3)p-(\sigma+7)\bar d
                   +2\sqrt{6(1-\bar d)(\sigma+1)p},
\]
\[
 Y_1=3(\sigma+1)q-(\sigma+5)\bar d
                   +2\sqrt{6(1-\bar d)(\sigma-1)q}.
\]
On $[1/3,2/5]$, both $Y_j/v^{3/2}>100$. To see this without a grid,
divide the positive terms by $v^{3/2}$. For $Y_0$ the two denominators
are $\sqrt p\,q^{3/2}$ and $p q^{3/2}$; for $Y_1$ they are
$p^{3/2}\sqrt q$ and $p^{3/2}q$. Their respective maxima on this interval
are smaller than
\[
                         63/200,\quad93/500,\quad49/250,\quad19/125.
\]
Differentiating their squares proves the endpoint locations (the second
has its maximum at $p=2/5$). Also $v^{3/2}>13/125$,
$\sqrt{6(1-\bar d)(\sigma+1)}>24/5$, and
$\sqrt{6(1-\bar d)(\sigma-1)}>69/20$.
Replacing $\sigma$ by $3.162277660$ in positive terms and by
$3.162277661$ in negative terms now proves both bounds by rational arithmetic.

Since $v\ge2/9$, $n\ge42$ and $\lambda\le8/25$, one has
$x=\lambda/(nv)\le6/175$. The elementary inequality
$(1+x)^{3/2}\le1+3x/2+3x^2/8$ implies
$(1+x)^{3/2}\le1+(31/20)x$. Hence the hypotheses of Theorem~\ref{g:main}
yield
\[
 24+100\lambda\le\frac{129}{4}
               \left(1+\frac{31}{20}\frac{3\lambda}{28}\right),
 \qquad\text{and therefore}\qquad\lambda<9/100.
\]
There has been no optimization over independent moment or noise parameters.

\section{Interpolation preserves the Bernoulli reserve}
Put $t=np+H\sqrt{v+s}$, $k=\lfloor t\rfloor$, and $u=t-k\in[0,1)$.
The support assumptions give the crossing theorem of v13, including
$na_0\le t<n(1-a_1)$. Its exact weighted identity implies
\begin{equation}\label{g:cross}
 \Pp(S_n\le H)\le(1-u)F_B(k)+uF_B(k+1)+\mathcal C(t).
\end{equation}
The nonnegative weighted crossing loss was discarded, not estimated by
an unrelated law. The identity uses closed lower crossing and strict upper
crossing events and remains valid at every atom.

Schulz's theorem, in its published $p\in[1/3,2/3]$ range, gives for every
real argument the strict Bernoulli reserve
\[
 \sqrt n\{F_B(x)-\Phi((x-np)/\sqrt{nv})\}-C_E\rho_B
 <-\frac{\phi_0Q(p)}{6\sqrt v}\le-\frac{4\phi_0}{n},
 \qquad \rho_B=(1-2v)/\sqrt v.
\]
Only the two integer arguments $k,k+1$ are used in \eqref{g:cross}.
This formula is not extended below $p=1/3$.

We bound the normal interpolation error with its sign. Write $w=t-np$,
and $G(x)=\Phi((x-np)/\sqrt{nv})$. The interpolation remainder satisfies
\[
 \sqrt n\{(1-u)G(k)+uG(k+1)-G(t)\}
 \le\frac{\phi_0}{2nv^{3/2}}u(1-u)(u-w)_+.
\]
Indeed $G''(x)\le\phi_0(np-x)_+/(nv)^{3/2}$ on the entire segment.
Since $v+s<1/4$, one has $|w|<23/25$ and
$c=(np-q-t)_+<8/25$. For $H\le0$, the normal change from variance $v$ to
$v+s$ is favorable. The interpolation error is bounded by
\[
 \frac{2/5}{nv^{3/2}}\left(\frac{23}{200}+\frac2{27}\right),
\]
using $u(1-u)\le1/4$ and $u^2(1-u)\le4/27$.
For $H\ge0$ the interpolation error is at most
$(2/5)(2/27)/(nv^{3/2})$, and the variance repair is at most
\[
 \frac{\phi_0ws}{2v^{3/2}}
             <\frac{(2/5)(23/25)(9/100)/2}{nv^{3/2}}.
\]
This uses $1-(1+x)^{-1/2}\le x/2$; here $c=0$.

The third-moment correction is favorable throughout this theorem.
The actual marginal expansion and inward cubic bounds give
\[
 (v+s)^{3/2}(\rho-\rho_B)
 \ge s\left\{3p-\frac32(1-2v)-\frac{17}{250}
                           -\frac{3(1-2v)s}{8v}\right\}>0
 \quad(s>0).
\]
Indeed $p\ge1/3$, $1-2v\le5/9$, $v\ge2/9$ and
$s<9/(100\cdot42)$ make the braces positive. At $s=0$ the correction
is exactly zero. This favorable sign is not asserted on the whole left
range below $1/3$.

Finally Lemma~\ref{g:curvature} applies with
\[
 T=3/20,\quad M=141/1000,\quad\beta=41/100,\quad
 \sqrt T<97/250,\quad v^{3/2}>13/125.
\]
Let
\[
 S_c=\frac{3}{80}+c\frac{97}{250}
          +\frac{27}{256}\frac{141}{1000}\left(\frac53+\frac{5c}4\right).
\]
For $H\le0$ and $H\ge0$, respectively, the total unfavorable errors
in \eqref{g:cross}, multiplied by $n$, are strictly smaller than
\[
 E_-:=\frac{(41/100)S_{8/25}+(2/5)(23/200+2/27)}{13/125},
\]
\[
 E_+:=\frac{(41/100)S_0+(2/5)(2/27)+(2/5)(23/25)(9/100)/2}{13/125}.
\]
Exact rational arithmetic gives
\[
              E_-,E_+<4\frac{19943}{50000}-\frac1{10}
                         <4\phi_0-\frac1{10}.
\]
Combine this with the strict Bernoulli reserve to prove Theorem~\ref{g:main}.
Both signs of $H$, every phase $u\in[0,1)$, zero noise and all threshold
atoms have been included.

\section*{Dependencies and limits}
The geometric inequality is proved in\par\noindent
\path{mathematics/simplification-v13/2026-10-28/geometry-noise.tex}.\par\noindent
The exact weighted crossing identity is in the sibling \path{crossing.tex}.
The external Bernoulli theorem is Theorem~1, printed page~1, of Schulz's
\emph{The optimal Berry--Esseen constant in the binomial case} (2016),
\href{https://d-nb.info/1197702695/34}{original thesis}.
Finite Stirling is the same explicit remainder bound used in the retained
count normalization. No Fourier, Poisson or count-generation program is
an input to the new finite arithmetic.

For acquired witnesses the actual central gap, the coarse geometric budget,
left support strips and raw threshold still require their existing
contact, sharp-predecessor and sign-specific acquisition. Applying the
statement to a reflected branch requires these inputs for the reflected
law. No global attainment, independent finite base through 41, comparison
below $p=1/3$, or full conjecture is proved here.
\end{document}
